% =====================================================================
%  figures/bivariable-cuali-cuanti.tex
%  Boxplot: 'gender' x 'biology_score'.
%  (Reemplaza a img/bivariable_cuali_cuanti.png)
%
%  Valores desde la tabla de SPSS (tables/bivariable-genero-biologia):
%    female: min 79, Q1 82, mediana 83,5, Q3 86, max 90
%    male:   min 79, Q1 83, mediana 85,   Q3 86, max 90
% =====================================================================
\begin{figure}[H]
\centering
\begin{tikzpicture}
\begin{axis}[
  boxplot/draw direction=y,
  title={Distribución de \Variable{biology\_score} según género},
  ylabel={Puntaje de biología},
  ymin=76, ymax=92,
  xtick={1,2},
  xticklabels={female, male},
  xlabel={Gender},
  grid=major,
]
  % female
  \addplot+[boxplot prepared={
      lower whisker=79, lower quartile=82, median=83.5,
      upper quartile=86, upper whisker=90},
      fill=azulclaro, draw=azul] coordinates {};
  % male
  \addplot+[boxplot prepared={
      lower whisker=79, lower quartile=83, median=85,
      upper quartile=86, upper whisker=90},
      fill=naranja!40, draw=naranja] coordinates {};
\end{axis}
\end{tikzpicture}
\caption{Relación entre género y puntaje en biología.}
\label{fig:bivariable_cuali_cuanti}
\end{figure}
